\documentclass[11pt,a4paper]{article}
\usepackage[utf8]{inputenc}
\usepackage{graphicx}
\usepackage[left=2.5cm,top=3cm,right=2.5cm,bottom=3cm,bindingoffset=0.5cm]{geometry}
\usepackage{macros/AEDLogica, macros/AEDEspecificacion, macros/AEDTADs}
\usepackage{caratula}


\titulo{Trabajo práctico}
\subtitulo{El subtitulo adecuado}

\fecha{\today}

\materia{Algoritmos y Estructuras de Datos}
\grupo{Java con Mates} 


\integrante{Ayala, Ignacio}{196/25}{ayalanacho11@gmail.com} 
\integrante{Della Bonzana, Lara}{563/25}{dellabonzanalara@gmail.com}
\integrante{Morales, Santiago }{481/25}{santumorales11@gmail.com}
\integrante{Vannelli, Tiziana}{945/25}{tizivannelli@gmail.com}  %


% Declaramos donde van a estar las figuras
% No es obligatorio, pero suele ser comodo
\graphicspath{{../static/}}

% Asi pueden escribir nuevos comandos. 
% Este por ejemplo asegura q los nombres 
% que figuren con una tipografia diferenciada  
\newcommand{\Tipo}[1]{\mathsf{#1}} 
% la sintaxis es \newcommand{\nombreDeLaMacro}[cantidadDeParametros]{Lo que va ser remplazado por el macro} 
\newcommand{\norm}[1]{\vert #1\vert}



\begin{document}

\maketitle

\section{Macros para escribir en LaTeX}

Para facilitar la escritura del trabajo práctico en LaTeX, les proveemos de varias macros que les permitirán escribir especificaciones en lógica de primer orden.
Por ejemplo, podrán hacer uso de comandos que definen los conectores lógicos: $\land, \lor, \implica$ y también los operadores de la lógica trivaluada $\yLuego$, $\oLuego$, $\implicaLuego$.

Además, disponen de comandos para cuantificadores: $\paraTodo{x}{\mathbb{Z}}{x \geq 0 \lor x < 0}$, $\existe{x}{\mathbb{Z}}{x \geq 0}$. A estos comandos se les puede agregar un \texttt{Largo} al final para que ocupen múltiples renglones.

\section{Ejemplos de uso de las macros}


\begin{tad}{EjemploDeTAD}


\obs{secuenciaDeCaracteres}{\seq{\cha}}
\\
\obs{conjuntoDeTuplas}{\conj{\tupla{nombre: \seq{\cha}, apellido: \seq{\cha}}}}
\\


\begin{proc}{crearTAD}
{
\In caracteres: \seq{\cha},
\In secuenciaDeNombres: \seq{\tupla{nombre: \seq{\cha}, apellido: \seq{\cha}}}
}{
\Tipo{EjemploDeTAD}
}
    \requiere{predicadoUnilinea(caracteres)}
    \asegura{res.secuenciaDeCaracteres = caracteres}
    \aseguraLargo{ \norm{res.conjuntoDeTuplas} = \norm{secuenciaDeNombres} \land \\predicadoMultilinea(secuenciaDeNombres, res.conjuntoDeTuplas) }
\end{proc}

\predLargo{predicadoMultilinea}{s: \seq{\cha}, conjunto: \conj{\tupla{\seq{\cha}, \seq{\cha}}}}{
    \paraTodo{i}{\Z}{0 \leq i < \norm{s} 
    \implicaLuego s[i] \in c}
}

\pred{predicadoUnilínea}{s: \seq{\cha}}{
    \norm{s} > 0
}

\aux{unAuxiliar}{k: \Z,l: \Z}{\Z}{\norm{k-l}}

\end{tad}

%Aca arranca
\newpage

Usuario ES $\struct{id: \Z, millasTotales: \Z, millasDisponibles: \Z, historial: \seq{} , categoria: \Tipo{Categoria}}$ \newline
Promocion ES $\struct{nombre: \cha, factor: \Z, estado: \bool}$ \newline
Categoria ES $\cha$ \newline

\begin{tad}{AirMiles}

\obs{usuarios}{\conj{\Tipo{Usuario}}} \newline
\obs{promociones}{\conj{\Tipo{Promocion}}} \newline
\obs{promocionesConTopeActivas}{\dict{\Tipo{Promocion}}{\tupla{\Z, \Z}}} \newline


\begin{proc}{crearSistema}
{}{\Tipo{AirMiles}}
    \requiere{True}
    \asegura{\existe{p}{\Tipo{Promocion}}{p \in res.promociones \land p.nombre = 'NoPromo' \land p.factor = 1 \land p.estado = True}} 
    \asegura{res.usuarios = \emptyset}
    \asegura{\# res.promociones =1}

\end{proc}



\begin{proc}{registrarUsuario}
{\Inout a: \Tipo{AirMiles}, \In id: \Z}{}
    \requiere{! esUsuarioRegistrado(id, a) \land a = A0}
    \asegura{\#a.usuarios = \#A0.usuarios + 1}
    \asegura{\paraTodo{u}{\Tipo{Usuario}}{esUsuarioRegistrado(u.id, A0) \implica esUsuarioRegistrado(u.id, a)}}
    \asegura{\existe{u}{\Tipo{Usuario}}{!esUsuarioRegistrado(u.id, A0) \land esUsuarioRegistrado(u.id, a) \land u.id = id \land u.millasDisponibles = 0 \land u.millasTotales = 0 \land u.historial = [] \land u.categoria = 'Bronce'}}
    \asegura{a.promociones = A0.promociones}
    
\end{proc}

\begin{proc}{registrarUsuarioConReferidor}
{\In idNuevoUsuario: \Z, \In referidor: \Tipo{Usuario}, \Inout a: \Tipo{AirMiles}}{}
    \requiere{esUsuarioRegistrado(referidor.id, a) \land a = A0 \land !esUsuarioRegistrado(idNuevoUsuario, a)}
   \aseguraLargo{\existe{u}{\Tipo{Usuario}}{esUsuarioRegistrado(u.id, a) \land u.id = idNuevoUsuario \land u.millasTotales = 500 \land u.millasDisponibles = 500 \land u.historial = [] \land u.categoria = Bronce}}
    \asegura{elRestoDeUsuariosSeMantienenIgual(referidor, A0, a)}
    \asegura{\#a.usuarios = \#A0.usuarios + 1}
    \aseguraLargo{\existe{u}{\Tipo{Usuario}}{esUsuarioRegistrado(u.id,a) \land u.id = referidor.id \land u.millasDisponibles = referidor.millasDisponibles + 500 \land u.millasTotales = referidor.millasTotales + 500 \land u.historial = referidor.historial ++ []\land u.categoria = queCategoriaEs(u.millasTotales,a)}}
    \asegura{a.promociones = A0.promociones}

\end{proc}

\begin{proc}{acumularMillas}
{\In usuario: \Tipo{Usuario}, \In distancia: \Z, \In nombrePromo: \cha, \Inout a:\Tipo{AirMiles}}{}
    \requiereLargo{esUsuarioRegistrado(usuario.id, a) \land a = A0 \land distancia \geq 0 \land esPromoActiva(nombrePromo, a)}
    \asegura{h}
\end{proc}


\begin{proc}{canjearMillas}
{\In u: \Tipo{Usuario}, \In cantMillas: \Z, \In premio: \Tipo{T}, \Inout a: \Tipo{AirMiles}}{}
    \requiereLargo{cantMillas > 0 \land esUsuarioRegistrado(u.id,a) \land a = A0 \land u.millasDisponibles \geq cantMillas \land u.millasTotales \geq u.millasDisponibles + cantMillas}
    \asegura{\#A0.usuarios = \#a.usuarios}
    \aseguraLargo{\existe{u_0}{\Tipo{Usuario}}{esUsuarioRegistrado(u_0.id,a) \land u_0.id = u.id \land u_0.millasDisponibles = u.millasDisponibles - cantMillas \land u.millasTotales = u_0.millasTotales \land u_0.historial = u.historial ++[] \land u_0.categria = u.categoria}}
    \asegura{A0.promociones = a.promociones}
    \asegura{elRestoDeUsuariosSeMantienenIgual(u, A0, a)}
\end{proc}

\begin{proc}{transferirMillas}
{\In cantMillas: \Z, \In usuarioOrigen: \Tipo{Usuario}, \In usuarioDestino: \Tipo{Usuario}, \Inout a: \Tipo{AirMiles}}{}
    \requiereLargo{esUsuarioRegistrado(usuarioDestino.id, a) \land esUsuarioRegistrado(usuarioOrigen.id, a) \land cantMillas \geq 0 \land usuarioOrigen.millasDisponibles \geq cantMillas \land usuarioDestino \neq usuarioOrigen \land a = A0 \land usuarioDestino.millasTotales \geq usuarioDestino.millasDisponibles + cantMillas}
    \asegura{\paraTodo{u}{\Tipo{Usuario}}{esUsuarioRegistrado(u.id, A0) \land u \neq usuarioOrigen \land u \neq usuarioDestino \implica esUsuarioRegistrado(u.id, a)}}
    \asegura{\#A0.usuarios = \#a.usuarios}
    \aseguraLargo{\existe{u_0}{\Tipo{Usuario}}{esUsuarioRegistrado(u_0.id, a) \land u_0.id = usuarioOrigen.id \land u_0.millasDisponibles = usuarioOrigen.millasDisponibles - cantMillas \land u_0.millasTotales = usuarioOrigen.millasTotales \land u_0.historial = usuarioOrigen.historial ++ [] \land u_0.categoria = usuarioOrigen.categoria}}
    \aseguraLargo{\existe{u_0}{\Tipo{Usuario}}{esUsuarioRegistrado(u_0,a) \land u_0.id = usuarioDestino.id \land u_0.millasDisponibles = usuarioDestino.millasDisponibles + cantMillas \land u_0.millasTotales = usuarioDestino.millasTotales \land u_0.historial = usuarioDestino.historial ++ [] \land u_0.categoria = usuarioDestino.categoria}}
\end{proc}

\begin{proc}{consultarSaldo}
{\In u: \Tipo{Usuario}, \In a: \Tipo{AirMiles}}{\Z}
    \requiere{esUsuarioRegistrado(u.id, a)}
    \asegura{res = u.millasDisponibles}
\end{proc}

\begin{proc}{consultarCategoria}
{\In u: \Tipo{Usuario}, \In a: \Tipo{AirMiles}}{\Tipo{Categoria}}
    \requiere{esUsuarioRegistrado(u.id, a)}
    \asegura{res = u.categoria}
\end{proc}

\begin{proc}{rankingMillas}
{tipos}{\Tipo{AirMiles}}
    \requiere{True}
    \asegura{algo}
\end{proc}

\begin{proc}{resumenTransacciones}
{tipos}{\Tipo{AirMiles}}
    \requiere{True}
    \asegura{algo}
\end{proc}

\begin{proc}{iniciarPromocion}
{\In nombre: \cha, \In factor: \ent, \Inout  a: AirMiles}{}
    \requiere{ factor \geq 1 \land a = A0}
    \asegura{\existe{p}{\Tipo{Promocion}}{p \in a.promociones \land p.nombre = nombre \land p.factor = factor \land p.estado = True}}
    \asegura{elRestoDePromocionesSeMantienenIgual(nombre, A0, a) \land elRestoDePromocionesSeMantienenIgual(nombre, a, A0)}
    \asegura{a.usuarios = A0.usuarios}
\end{proc}
\newline Observacion: Usamos el predicado dos veces pero con distinto orden para afirmar que toda promocion esta en un sistema si y solo si esta en el otro \newline
\begin{proc}{iniciarPromocionConTope}
{\In nombre: \cha, \In factor: \ent, \In tope: \ent, \Inout a: \Tipo{AirMiles}}{\Tipo{AirMiles}}
    \requiere{a = A0 \land factor \geq 0 \land tope \geq 0}

    \asegura{elRestoDePromocionesSeMantienenIgual(nombre, A0, a) \land elRestoDePromocionesSeMantienenIgual(nombre, a, A0)}
\end{proc}

\begin{proc}{finalizarPromocion}
{\In nombre: \cha, \Inout a: \Tipo{AirMiles}}{}
    \requiere{ esPromoActiva(nombre, a) \land a=A0 \land nombre \neq NoPromo}
    
    \asegura{\existe{p}{\Tipo{Promocion}}{p.nombre = nombre \land p.estado = False \land p \in a.promociones }}
    \asegura{elRestoDePromocionesSeMantienenIgual(nombre, A0, a)}
    \asegura{\#a.promociones = \#A0.promociones}
    \asegura{a.usuarios = A0.usuarios}
\end{proc}

\begin{proc}{elQueTieneMasMillas}
{\In a: \Tipo{AirMiles}}{\Z}
    \requiere{True}
    \asegura{\existe{u}{\Tipo{Usuario}}{esUsuarioRegistrado(u.id, a) \land u.id = res \land tieneMaximasMillas(u, a)}}
\end{proc}

\begin{proc}{elMasCanjeador}
{tipos}{\Tipo{AirMiles}}
    \requiere{True}
    \asegura{algo}
\end{proc}

\begin{proc}{paresTransferidoresExclusivos}
{tipos}{\Tipo{AirMiles}}
    \requiere{True}
    \asegura{algo}
\end{proc}

\aux{queCategoriaEs}{millas: \Z, a: \Tipo{AirMiles}}{\Tipo{Categoria}}{
     \IfThenElse{millas \leq 10000}{Bronce}{\newline \IfThenElse{millas\leq 50000}{Plata}{IfThenElse\{millas\leq 1000000, Oro, Platino\}}}} \newline

\aux{factorDeCategoria}{categoria: \Tipo{Categoria}}{\Z}{\IfThenElse{categoria = Bronce}{1}{\IfThenElse{categoria = Plata}{1.5}{\IfThenElse{categoria = Oro}{2}{0}}}}

\aux{calculoMillas}{distancia: \Z, factorCategoria: \Z, factorPromocion: \Z}{\Z}{distancia * factorCategoria * factorPromocion}

\pred{tieneMaximasMillas}{u: \Tipo{Usuario}, a: \Tipo{AirMiles}}{\paraTodo{u_0}{\Tipo{Usuario}}{u_0 \in a.usuarios \land u_0 \neq u \implica u.millasTotales \geq u_0.millasTotales}}\newline
\pred{esUsuarioRegistrado}{id: \Z, a: \Tipo{AirMiles}}{
    \existe{u_0}{\Tipo{Usuario}}{u_0 \in a.usuarios \yLuego u_0.id = id}
} \newline
\pred{esPromoActiva}{nombre: \cha, a:\Tipo{AirMiles}}{
    (\existe{p}{\Tipo{Promocion}}{p \in a.promociones \yLuego p.nombre = nombre \land p.estado = True})
} \newline

\pred{elRestoDeUsuariosSeMantienenIgual}{usuario: \Tipo{Usuario}, A0: \Tipo{AirMiles}, a:\Tipo{AirMiles}}{\paraTodo{u}{\Tipo{Usuario}}{u \in A0.usuarios \land u \neq usuario \implicaLuego u \in a.usuarios}} \newline
\pred{elRestoDePromocionesSeMantienenIgual}{nombrePromocion: \cha, A0: \Tipo{AirMiles, a: \Tipo{AirMiles}}}{\paraTodo{p}{\Tipo{Promocion}}{p \in A0.promociones \land p.nombre \neq nombrePromocion \implica p \in a.promociones}} \newline
\pred{esPromocionDelSistema}{promo: \Tipo{Promocion} a: \Tipo{AirMiles}}{p \in a.promociones} \newline
\pred{tienenMismoId}{u_0: \Tipo{Usuario}, u_1: \Tipo{Usuario}}{u_0.id = u_1.id} \newline
\pred{tienenMismasMillasDisponibles}{u_0: \Tipo{Usuario}, u_1: \Tipo{Usuario}}{u_0.millasDisponibles = u_1.millasDisponibles} \newline
\pred{tienenMismasMillasTotales}{u_0: \Tipo{Usuario}, u_1: \Tipo{Usuario}}{u_0.millasTotales = u_1.Totales} \newline
\pred{tienenMismaCategoria}{u_0: \Tipo{Usuario}, u_1: \Tipo{Usuario}}{u_0.categoria= u_1.categoria} \newline
\pred{tienenMismoHistorial}{u_0: \Tipo{Usuario}, u_1: \Tipo{Usuario}}{u_0.historial = u_1.historial} \newline
\pred{tieneSaldoSuficiente}{u: \Tipo{Usuario}, cantMillas : \Z}{u.millasDisponibles \geq cantMillas} \newline
\pred{NoSuperaMillasTotales}{u: \Tipo{Usuario}, cantMillas: \Z}{u.millasTotales \geq u.millasDisponibles + cantMillas} \newline
\end{tad}

\end{document}
